\section{Limitations}\label{sec:limits}
\begin{itemize}
\item \textbf{Conditional reference.} Classical LMMSE optimality assumes
specified covariance, noise and operator. A generic spectrum is not a
universal bound; individual-image periodograms are oracle references.
\item \textbf{Band approximation.} Scalar decoupling need not hold for
aliased or rotated observations. Band averaging has no uniform error
direction. Windows, cropping, clipping and solver floors create additional
departures from exact theory.
\item \textbf{Network mechanism.} G1 and L2 establish empirical ordering,
not a gain-mismatch mechanism or function equivalence. L2 function
similarity passes in 0/4 groups; superiority over the simple score passes only
2/4. No scalar-scale-only retraining arm has been evaluated. Exploratory
same-observation controls in Appendix~\ref{app:controls} further retain
opposing kernel/FF orderings; neither finite Fourier ridge nor a model's
own posterior variance establishes the trained network mechanism.
\item \textbf{Failed tests and sensitivity.} F1, E1 switching, D11 winner
selection and L1 strict collapse retain their failures. E3 similarity
passes 10/12 at .03 and 8/12 at .02; its .1/.8 reference-regime thresholds
remain descriptive pending the sensitivity audit.
\item \textbf{Replication scope.} G1/L2 cover two additional fixed seeds
and two pairs. G1 uses HR256 and 73 images ($\times3$ means 256/85), while
L2 uses HR240 and 74. Their pooled rows cannot be interchanged. G1 saves
per-image recovery but not signal/error energies, limiting independent
reconstruction of its pooled statistics.
\item \textbf{Protocol and numerical differences.} E2 windowing changed
after a one-image pilot. E3 has no injected observation noise although its
reference uses $\sigma=.002$ and its GP solver $\nu=10^{-5}$. E1's native
forward geometry differs from its inverse discretisation, and clipped
outputs are not fixed linear estimators. D11's revised bootstrap covers
fixed images/noise draws and does not measure training-seed variability.
\item \textbf{Guidance generalisation.} G2 covers 12 fixed images, three
noise settings with fixed draws, one operator and one 32-image calibration.
Zero observed severe degradations does not certify a population rate below
5\%. Analysis chooses the same prior in all cases, so its improvement
over the default does not establish target-specific adaptation. It trails
CV in accuracy and cost. Four-probe variability, calibration mismatch and
unobserved spectral components remain limitations; no INR architecture
selection is established. Oracle candidates were not all saved as pixels.
\item \textbf{Action-class dependence.} Unidentifiable spectra can still
permit useful mixture reconstruction. Appendix~\ref{app:ambiguity}
distinguishes selecting a fixed candidate from arbitrary estimators;
its elementary two-point separation is classical illustrative material,
not a general sharp sampling theorem or a nonlinear-network guarantee.
\item \textbf{EEG evidence.} EEG currently contributes 200 calibration segments;
the 760-segment test result has not been recovered and is excluded from
confirmatory claims in this revision.
\end{itemize}

\section{Conclusion}
Observation-conditioned linear references clarify spectral mismatch costs
and their asymmetry. Exact scalar and matrix identities provide the
foundation; band approximations and neural interpretations require
separate evidence. Completed experiments support mismatch-cost ranking
and repeated readout-loss ordering while preserving failures of the earlier
winner diagnostic, function equivalence and a proposed gain mechanism. Practical
guidance must specify the reference spectrum and operator and demonstrate
selection regret against competing baselines. A completed prospective test
shows that calibrated analysis can improve a fixed GP default, but its
constant choice, higher regret and higher cost than CV leave adaptive and
computational advantages unsupported. The formulas are useful conditional
diagnostics; deployment benefits require direct comparative evidence.
